\section{Method}
\label{sec:method}

\begin{figure*}[t]
\centering
% =========================================================== (a) what A* tracks
\begin{minipage}[b]{0.56\textwidth}\centering
\begin{tikzpicture}
\begin{axis}[width=\linewidth, height=4.4cm,
  xmin=-0.55, xmax=5.55, ymin=-1.75, ymax=3.75,
  ylabel={cap angle $\alpha$ [rad]}, ylabel style={yshift=-4pt},
  xtick={0,...,5},
  xticklabels={start, enter, stroke 1, return, stroke 2, release},
  x tick label style={font=\scriptsize},
  ytick={-1,0,1,2,3}, font=\scriptsize,
  grid=major, grid style={gray!15}, clip=false]
  \addplot[dashed, red!70!black, domain=-0.55:5.55, samples=2] {3};
  \node[font=\scriptsize, red!70!black, anchor=north west] at (axis cs:-0.5,3.0) {$\agoal$};
  \foreach \i/\lo/\hi in {0/-1/-1, 1/-1/-1, 2/-1/1, 3/-1/1, 4/-1/3, 5/0.35/3} {
    \edef\temp{\noexpand\fill[blue!18, draw=blue!55, rounded corners=1pt]
      (axis cs:\i-0.24,\lo-0.07) rectangle (axis cs:\i+0.24,\hi+0.07);}
    \temp
  }
  \addplot[thick, black!70, densely dotted] coordinates {(0,-1) (1,-1)};
  \addplot[very thick, green!45!black] coordinates {(1,-1) (2,1)};
  \addplot[thick, orange!85!black, dashed] coordinates {(2,1) (3,1)};
  \addplot[very thick, green!45!black] coordinates {(3,1) (4,3)};
  \addplot[thick, black!70, densely dotted] coordinates {(4,3) (5,3)};
  \addplot[only marks, mark=*, mark size=1.3pt, black] coordinates {(0,-1) (1,-1) (2,1) (3,1) (4,3) (5,3)};
  \node[font=\scriptsize, green!35!black] at (axis cs:3,-1.5)
    {each stroke widens the reachable interval};
  \draw[-{Stealth[length=1.6mm]}, blue!60!black] (axis cs:4.75,-0.55) -- (axis cs:4.3,-0.2);
  \node[font=\scriptsize, blue!60!black, anchor=west, align=left] at (axis cs:4.6,-0.9) {reachable\\[-2pt]cap angles};
  \node[font=\scriptsize, black!70, anchor=south west] at (axis cs:4.55,3.05) {goal \checkmark};
  \fill[blue!12, draw=blue!50] (axis cs:-0.3,-2.95) rectangle (axis cs:0.3,-2.55);
  \node[font=\scriptsize] at (axis cs:0,-2.75) {$F_2$};
  \fill[green!15, draw=green!45!black] (axis cs:0.7,-2.95) rectangle (axis cs:4.3,-2.55);
  \node[font=\scriptsize, green!25!black] at (axis cs:2.5,-2.75) {transition set $T_6=G\cap F_6$};
  \fill[blue!12, draw=blue!50] (axis cs:4.7,-2.95) rectangle (axis cs:5.3,-2.55);
  \node[font=\scriptsize] at (axis cs:5,-2.75) {$F_0$};
  \node[font=\scriptsize, anchor=east, gray] at (axis cs:-0.38,-2.75) {set};
\end{axis}
\end{tikzpicture}\\[-2pt]
{\footnotesize (a) What the search tracks}
\end{minipage}\hfill
% =========================================================== (b) wrist-cap plane
\begin{minipage}[b]{0.40\textwidth}\centering
\begin{tikzpicture}
\begin{axis}[width=\linewidth, height=4.4cm,
  xmin=-1.15, xmax=1.15, ymin=-1.75, ymax=3.75,
  xlabel={wrist angle $w$ [rad]}, ylabel={cap angle $\alpha$ [rad]},
  ylabel style={yshift=-4pt}, xlabel style={yshift=3pt},
  ytick={-1,0,1,2,3}, xtick={-1,0,1}, font=\scriptsize,
  grid=major, grid style={gray!15}]
  \addplot[thick, black!70, densely dotted, -{Stealth[length=1.6mm]}] coordinates {(0.04,-1) (-0.98,-1)};
  \addplot[very thick, green!45!black, -{Stealth[length=1.8mm]}] coordinates {(-1,-1) (1,1)};
  \addplot[thick, orange!85!black, dashed, -{Stealth[length=1.6mm]}] coordinates {(1,1) (-0.98,1)};
  \addplot[very thick, green!45!black, -{Stealth[length=1.8mm]}] coordinates {(-1,1) (1,3)};
  \addplot[thick, black!70, densely dotted, -{Stealth[length=1.6mm]}] coordinates {(1,3) (0.06,3)};
  \addplot[only marks, mark=*, mark size=1.4pt] coordinates {(0.04,-1) (0.04,3)};
  \node[font=\scriptsize, anchor=north] at (axis cs:0.04,-1.08) {start};
  \node[font=\scriptsize, anchor=south] at (axis cs:0.04,3.05) {goal};
  \node[font=\scriptsize, green!35!black, rotate=36] at (axis cs:0.12,-0.35) {stroke 1};
  \node[font=\scriptsize, green!35!black, rotate=36] at (axis cs:0.12,1.65) {stroke 2};
  \node[font=\scriptsize, orange!80!black, anchor=south] at (axis cs:-0.45,1.02) {return};
  \node[font=\scriptsize, anchor=west, align=left] at (axis cs:-1.1,2.3) {$\Delta\alpha=\Delta w$\\[-2pt]while grasping};
\end{axis}
\end{tikzpicture}\\[-2pt]
{\footnotesize (b) The same plan in the wrist--cap plane}
\end{minipage}
\caption{The reachability A* on the running example, task T2 of
Section~\ref{sec:experiments}, query Q0: turn the cap from $\alpha=-1$ to
$3$\,rad and bring the arm back to its start.
(a)~Each column is one node on the returned path: the set the robot is in
(strip) and the interval of cap angles reachable there (bar). The black line
is the cap angle along the final plan. Entering the transition set $T_6$
changes nothing; each stroke widens the interval by up to the wrist range of
$T_6$ (2\,rad); returns and transit keep it. Leaving $T_6$ towards $F_0$ needs
some further wrist motion, which raises the foot of the last bar. The search
stops at the first node whose set contains the goal arm configuration and
whose interval contains $\agoal$. (b)~While grasping, the cap turns as the
wrist moves forward; returns move the wrist back with the cap fixed. Two
strokes are the minimum, since each turns the cap by at most 2\,rad.
\todo{regenerate once the transition sets are validated (E5).}}
\label{fig:reach}
\end{figure*}

The planner has an offline stage, which depends only on the robot, the object
and the scene, and an online stage that is run for each query. Offline, it
builds the free sets (Section~\ref{sec:method:free}) and the grasp set
(Section~\ref{sec:method:grasp}), and intersects them to obtain validated
transition sets and the graph over them, using linear programs alone
(Section~\ref{sec:method:graph}). Online, the reachability A* chooses the
sequence of sets and grasps (Section~\ref{sec:method:search}), and a second
stage turns it into motions (Section~\ref{sec:method:recovery}).
New goals and start configurations reuse the offline stage unchanged.
Fig.~\ref{fig:reach} illustrates the online stage on the task of
Section~\ref{sec:experiments}, which we use throughout as a running example.

\subsection{Free sets}
\label{sec:method:free}
We inner-approximate $\CSfree$ by polytopes $P_j=\{q:A_jq\le b_j\}$,
$j=1,\dots,J$, computed with clique covers of visibility
graphs~\cite{werner2024approximating} and IRIS-NP2
regions~\cite{werner2024faster} in Drake~\cite{tedrake2019drake}. By
Assumption~\ref{ass:task}(iii), the object coordinate does not affect
collisions. We therefore replace each region by its slice at $\alpha=0$,
$F_j=\{x:(x,0)\in P_j\}$. This is sound: if $(x,0)$ is collision-free, so is
$(x,\alpha)$ for every $\alpha$. The robot can then occupy any $x\in F_j$ at
any object angle, and the planner tracks $\alpha$ separately.

\subsection{Grasp set}
\label{sec:method:grasp}
A grasp fixes the pose of the hand frame $E$ relative to the object frame $C$,
a nonlinear function of the arm joints. We write the resulting conditions as
\begin{equation}
  h(x)=0,\qquad g(x)\le 0,
  \label{eq:grasp_hg}
\end{equation}
where $h$ collects the equalities a valid grasp must satisfy, typically the
components of ${}^{C}p^{E}(x)$ and of the hand axes that the grasp pins down,
and $g$ the inequalities, such as tolerance bands on the remaining components,
an approach direction and a finger opening.
Section~\ref{sec:experiments} instantiates them for our tasks. The set
$\{h=0,\,g\le0\}$ is thin and curved, so it cannot be written as a polytope.
We approximate it locally, using two ideas from NLP
sampling~\cite{toussaint2024nlp}. First, a seed grasp $x_0$ is found by
Gauss--Newton minimisation of the constraint violation
$\|[h(x);\max(g(x),0)]\|^2$ from a random start. Second, the directions that
preserve the grasp to first order are the null space of the Jacobian $J_h$ of
the equality constraints, which NLP sampling uses to move along constraint
manifolds. The grasp set keeps the configurations that stay within $\epsilon$
of this null space, around $x_0$:
\begin{equation}
\begin{aligned}
G=\{x:\ & |h_0+J_h(x-x_0)|\le\epsilon,\ \ g_0+J_g(x-x_0)\le\epsilon,\\
        & x^{G}_{\min}\le x\le x^{G}_{\max},\ \ |x_j-x_{0,j}|\le\rho\ \ \forall j\in\mathcal{A}\},
\end{aligned}
\label{eq:graspPoly}
\end{equation}
where $h_0,J_h,g_0,J_g$ are evaluated at $x_0$, $\epsilon$ is a tolerance, and
$\mathcal{A}$ are the joints other than the wrist.

Two properties shape this set. First, whenever a joint's motion leaves $h$ and
$g$ unchanged, its columns in $J_h$ and $J_g$ vanish and $G$ spans that joint's
entire range. This is what allows a stroke to run inside a single set, and it
is exact rather than approximate: it holds for the wrist of any arm whose last
axis passes through the origin of $E$ along the approach direction, since
rotating about it moves neither the position nor the approach axis of the hand.
Second, $J_h$ leaves a null space of dimension $n_r-\mathrm{rank}(J_h)$ that the
constraints of~\eqref{eq:graspPoly} bound only through the joint limits. Far
from $x_0$ the linearisation is meaningless, so the trust region $\rho$ keeps
the set where the linearisation is valid; Section~\ref{sec:results} quantifies
what happens without it. Seeds on other parts of the grasp manifold give
further grasp sets that the search treats in the same way. We use a single
seed in this paper (see Section~\ref{sec:discussion}).

\subsection{Transition sets and the graph}
\label{sec:method:graph}
For every free set with $F_k\cap G\ne\emptyset$, we form the transition set
$T_k=G\cap F_k$. Two LPs give its extreme wrist values,
$w_k^{+}=\max\{w:x\in T_k\}$ and $w_k^{-}=\min\{w:x\in T_k\}$, together with
points $x_k^{+},x_k^{-}\in T_k$ attaining them. A stroke inside $T_k$ turns the
object by at most $L_k=w_k^+-w_k^-$, and we discard sets with $L_k<\Lmin$.
Since $G$ is thin in every direction but the wrist, $T_k$ is a sliver along the
wrist axis, and by convexity the whole segment $[x_k^-,x_k^+]$ lies inside it:
the robot can sweep the entire wrist range of $T_k$ without leaving it.

\paragraph{Validation} The free sets are collision-free only in a volume
sense: IRIS-NP2 bounds the fraction of each region's volume in
collision~\cite{werner2024faster}. A transition set has almost no volume, so it
can lie entirely in collision without violating that bound. We therefore
validate each $T_k$ on its own. Because it is a sliver, this is cheap: we test
densely spaced configurations along $[x_k^-,x_k^+]$ with a collision checker,
and keep only the wrist sub-intervals on which every test passes. Each of these is $T_k$ intersected with a bound on
$w$, and is therefore still a polytope. \todo{implement (E5); state the check
resolution; this is a dense check, not a certificate. Certifying these thin
sets with C-IRIS~\cite{dai2024certified} is an alternative.}

\paragraph{Graph} The graph has two kinds of edges. A \emph{transit} edge
joins $F_i$ and $F_j$ if they intersect, and its door point is the Chebyshev
centre of $F_i\cap F_j$. A \emph{switch} edge joins $F_i$ and $T_k$ if they
intersect. For it, we store the wrist range $[a_{ik}^-,a_{ik}^+]$ of
$F_i\cap T_k$ and the points attaining its ends. Each of these quantities is
an LP feasibility check, a Chebyshev centre or an extreme point. No integer
variables or vertex enumeration are needed, and the number of LPs is
$O(J^2+JK)$ for $J$ free sets and $K$ transition sets.

\subsection{Reachability A*}
\label{sec:method:search}
We use A* in the same form and with the same syntax
as~\cite{griffin2019footstep,wang2026cassr} (Algorithm~\ref{alg:astar}).
Only the node structure, the expansion, the check for nodes already expanded,
the goal test and the costs are specific to our problem. They are highlighted
in blue and described below.

\begin{algorithm}[t]
\caption{A* search; the parts specific to our problem are in blue}
\label{alg:astar}
\begin{algorithmic}[1]
\State \hl{addStartNodesToQueue()} \label{ln:start} \Comment{\hl{Sec.~\ref{sec:astar:astar}}}
\While{hasNodesToCheck()}
  \State nodeToExpand = getCheapestNode() \label{ln:cheapest}
  \If{\hl{nodeAlreadyExpanded()}} \label{ln:skip} \Comment{\hl{Sec.~\ref{sec:astar:dominance}}}
    \State skipToNextNodeInQueue()
  \EndIf
  \If{\hl{hasReachedTheGoal()} or hasTimedOut()} \label{ln:goal}
    \State stopSearch()
  \EndIf
  \State childNodesToCheck = \hl{expandNode()} \label{ln:expand} \Comment{\hl{Alg.~\ref{alg:expand}}}
  \For{childNode $\in$ childNodesToCheck} \label{ln:forchild}
    \If{\hl{nodeAlreadyExpanded(childNode)}} \label{ln:skipchild}
      \State \textbf{continue} \Comment{skip node}
    \EndIf
    \State costOfEdge = \hl{getCostOfEdge()} \Comment{\hl{Sec.~\ref{sec:astar:cost}}}
    \State addEdgeToGraph(childNode, costOfEdge)
    \State costToGoal = \hl{estimateCostToGoal()} \Comment{\hl{Eq.~\eqref{eq:heuristic}}}
    \State nodeCost = costOfPath(childNode) + costToGoal \label{ln:nodecost}
    \State addNodeToQueue(childNode, nodeCost) \label{ln:addqueue}
  \EndFor
\EndWhile
\State plan = getBestPathToEndNode() \label{ln:plan} \Comment{\hl{\textsc{none} if queue empty}}
\end{algorithmic}
\end{algorithm}

\subsubsection{The A* algorithm}
\label{sec:astar:astar}
A* keeps a priority queue ordered by the cost of the path to each node plus
an estimate of the cost to the goal (line~\ref{ln:nodecost}). It starts with
one node per free set containing $\xinit$, with the interval $\{\ainit\}$
(line~\ref{ln:start}). At each iteration, it takes the cheapest node
(line~\ref{ln:cheapest}) and skips it if an equivalent or better node was
already expanded (line~\ref{ln:skip}). If the node is a goal
(line~\ref{ln:goal}), the plan is read by following parents back to the start
(line~\ref{ln:plan}). Otherwise, the node is expanded (line~\ref{ln:expand}),
and its children are costed and queued
(lines~\ref{ln:forchild}--\ref{ln:addqueue}). If the queue empties first, no
plan exists within the sets.

\begin{nodestruct}[t]
\caption{Node}
\label{struct:node}
\begin{algorithmic}[0]
\State parent: \textsc{Node}* \Comment{parent node}
\State setId: \textsc{Int} \Comment{free set $F_i$ or transition set $T_k$}
\State \hl{interval: $[\alo,\ahi]$} \Comment{\hl{reachable object configurations}}
\State \hl{phase: \textsc{Enum}} \Comment{\hl{transition sets: entered, holding, released}}
\State \hl{entryRange: $[a^-,a^+]$} \Comment{\hl{transition sets: wrist range at entry}}
\end{algorithmic}
\end{nodestruct}

\subsubsection{Node structure}
\label{sec:astar:node}
As in CASSR, a node does not represent a configuration but a set reachable
from the start (Struct~\ref{struct:node}). A node in a free set $F_i$ with
interval $[\alo,\ahi]$ states that every $x\in F_i$ is reachable together with
every $\alpha\in[\alo,\ahi]$; since $F_i$ is convex and $\alpha$ is constant
during transit, a straight line joins any two of its points without leaving
it. A node
in a transition set $T_k$ also records the wrist range $[a^-,a^+]$ through
which $T_k$ was entered, and a phase: \emph{entered} (not yet grasping),
\emph{holding} (after a stroke) or \emph{released} (after a return, with the
wrist at $w_k^-$). The interval carries all the information about the path so
far: reaching $F_3$ before any grasp and after a stroke of 2\,rad gives the
intervals $\{-1\}$ and $[-1,1]$, two different nodes.

\begin{algorithm}[t]
\caption{expandNode, for a node in set $S$ with interval $[\alo,\ahi]$, phase
$\phi$ and entry range $[a^-,a^+]$}
\label{alg:expand}
\begin{algorithmic}[1]
\State nodeList = []
\If{$S=F_i$ is a free set}
  \For{all $F_j$ with $F_i\cap F_j\neq\emptyset$} \Comment{transit}
    \State nodeList.add($F_j$, $[\alo,\ahi]$)
  \EndFor
  \For{all $T_k$ with $F_i\cap T_k\neq\emptyset$} \Comment{enter}
    \State nodeList.add($T_k$, $[\alo,\ahi]$, entered, $[a_{ik}^-,a_{ik}^+]$)
  \EndFor
\Else \Comment{$S=T_k$ is a transition set}
  \If{$\phi\in\{$entered, released$\}$} \Comment{stroke}
    \State $c\gets w_k^+-a^-$ if $\phi=$ entered, else $c\gets L_k$
    \State $\beta\gets\min(\ahi+c,\,\agoal,\,\alpha_{\max})$
    \State nodeList.add($T_k$, $[\alo,\beta]$, holding)
  \EndIf
  \If{$\phi=$ holding} \Comment{return}
    \State nodeList.add($T_k$, $[\alo,\ahi]$, released)
  \EndIf
  \If{$\phi\in\{$entered, holding$\}$} \Comment{release or leave}
    \For{all $F_j$ with $F_j\cap T_k\neq\emptyset$}
      \State $d\gets\max(0,a_{jk}^- - a^+)$ \Comment{forced turning}
      \If{$\phi=$ holding or $d=0$}
        \State nodeList.add($F_j$, $[\alo+d,\ahi]$)
      \EndIf
    \EndFor
  \EndIf
\EndIf
\State \Return nodeList
\end{algorithmic}
\end{algorithm}

\subsubsection{expandNode()}
\label{sec:astar:expand}
Algorithm~\ref{alg:expand} lists the children of a node. A stroke turns the
object by any amount up to $c$, the wrist travel left in $T_k$ ($w_k^+-a^-$
after entering, $L_k$ after a return), so it raises $\ahi$ and keeps $\alo$.
Since the object can only be grasped inside transition sets, strokes are the
only moves that widen the interval. Leaving $T_k$ towards $F_j$ while holding
the object is a release, and it can force some turning: the wrist only advances
while grasping, so reaching the exit range $[a_{jk}^-,a_{jk}^+]$ from the entry
range turns the object by at least $d$, which raises $\alo$. Without holding,
the robot may only leave if $d=0$. Children with an empty interval are
discarded, and children in the same $T_k$ keep the entry range of their parent.
\tocheck{$d$ trims only the least-turned end of the interval, and
measuring it from $a^+$ is exact there: the least-turned configuration that can
leave is the one that entered as far forward as possible and turned just
enough. The open question is physical. A return winds the wrist back without
the object, which assumes that letting go frees the wrist; under that same
assumption no exit forces any turning, $d$ is always $0$, and leaving from
\emph{entered} needs no condition. The rule as written only makes sense if the
object has to be carried to the doorway, which contradicts the return. Decide
which model holds before reimplementing.}

\subsubsection{nodeAlreadyExpanded()}
\label{sec:astar:dominance}
A node is skipped if an expanded node with the same key has an interval that
contains its interval, at no higher cost. The key is the set for free nodes,
and the set, phase and entry range for transition nodes. The check is applied
when a node is taken from the queue (line~\ref{ln:skip}) and before a child is added to it
(line~\ref{ln:skipchild}). Because the key and the interval describe everything that can still
be reached from a node, this check is exact: a skipped node can reach nothing
that the expanded one cannot reach at the same or lower cost. In CASSR, two
nodes are instead treated as the same if their foothold sets lie within an
empirical 2\,cm of each other~\cite{wang2026cassr}. Unlike CASSR, we also allow
the search to return to a set it has already left, since the same set can be
revisited with a different interval, and dominance removes the revisits that
cannot help.

\subsubsection{hasReachedTheGoal()}
\label{sec:astar:goal}
A node is a goal if it lies in a free set $F_i$ with $\xgoal\in F_i$ and
$\agoal\in[\alo,\ahi]$.

\subsubsection{Cost computation}
\label{sec:astar:cost}
getCostOfEdge() returns 1 for a stroke, a return or a release, and $\mu\ll1$
($\mu=10^{-3}$) for every other move. Every stroke is followed by exactly one
return or release, so the cost of a plan is twice its number of grasps, plus
$\mu$ times its number of set changes. As long as a plan makes fewer than
$1/\mu$ set changes, minimising this cost minimises the number of grasps
first and the set changes second. Releasing must cost 1, or leaving and
re-entering a transition set would be a cheaper substitute for a return. \tocheck{the code also allows a return from the
\emph{entered} phase at cost 1, before any grasp; such a move should cost
$\mu$ for the cost to be exactly twice the number of grasps.}
estimateCostToGoal() returns, with $\Lmax=\max_k L_k$,
\begin{equation}
  h(\nu)=2\Big\lceil\frac{\agoal-\ahi}{\Lmax}\Big\rceil+\mathbf{1}[\phi=\text{holding}].
  \label{eq:heuristic}
\end{equation}
It is admissible: each remaining stroke turns the object by at most $\Lmax$
and is followed by a return or the final release, and a node that is holding
the object must still release it. It is also consistent, since no move lowers
$h$ by more than its cost. CASSR instead scales its distance-based heuristic,
which makes it inadmissible and turns the search into a weighted
A*~\cite{wang2026cassr}. With an admissible heuristic and exact dominance, the
first goal node expanded has the minimum cost.

\subsection{From sets to motions}
\label{sec:method:recovery}
The search fixes the sets and the grasp actions, e.g., \emph{enter, stroke,
return, stroke, release}. As in CASSR, a second stage computes the
configurations, here without a solver. A backward pass from $\agoal$ decides
how much each visit to a transition set turns the object; it ends exactly at
$\ainit$ because the search proved $\agoal$ reachable. Waypoints are placed at
door points for transit, in $F_i\cap T_k$ for entering and leaving, and on
$[x_k^-,x_k^+]$ for strokes and returns. Each motion is a straight line
between two waypoints of one set, so by convexity it lies in that set; we
check this and the coupling for every motion. Path length is not optimised: a
convex program over the chosen sets, e.g., GCS~\cite{marcucci2023motion},
could shorten the path. \todo{implement, or keep as future work.}

\subsection{Properties}
\label{sec:method:properties}
\begin{proposition}[informal]
\label{prop:guarantees}
Relative to the free sets $\{F_j\}$ and the transition sets $\{T_k\}$:
(i) every motion returned is a straight line inside one set that satisfies
the coupling of Assumption~\ref{ass:task}, so the path is collision-free
whenever the sets are;
(ii) the search terminates. If a plan exists whose transit motions lie in the
free sets and whose strokes lie in the transition sets, it returns one with
the minimum number of grasps, with ties broken by the number of set changes.
Otherwise, it returns \textsc{none};
(iii) the second stage succeeds on every sequence the search returns.
\end{proposition}
\noindent\emph{Sketch.} (i) follows from convexity of the recorded sets and
from the coupling check applied to every motion.
(ii) Every interval end the search produces is built from the finitely many
wrist values the LPs of Section~\ref{sec:method:graph} return, clipped at
$\agoal$, so only finitely many nodes exist per key and the search
terminates. With a
consistent heuristic and exact dominance, A* expands a goal node of minimum
cost first~\cite{hart1968astar}. Entering at the lowest and leaving at the highest wrist value of a switch edge
maximises the travel available, so restricting entries and exits to these
points loses no plan. \tocheck{the two choices are made by different paths, one realising each
end of the interval, which is what a reachable set allows; say so explicitly.}
(iii) The backward pass mirrors the interval updates of
Algorithm~\ref{alg:expand}. \tocheck{write the full proof; (iii) is asserted in our
implementation but not yet proven. Overlapping transition sets would need
chain-aware placement, so each set is its own segment for now.}
\hfill$\square$


\subsection{Relation to a mixed-integer formulation}
\label{sec:method:miqp}
The same problem can be posed as a mixed-integer program, in the way convex
decompositions are normally used for contact
planning~\cite{deits2014footstep}: for a fixed number of grasps $\ell$, choose
keyframes $\qinit,x_1,\dots,x_{2\ell},\qgoal$, let a binary $z_{ij}$ indicate
that keyframe $i$ lies in set $j$, and increase $\ell$ until the program is
feasible. Membership is an indicator constraint, usually in big-M form,
$A_jx_i\le b_j+M(1-z_{ij})$. Two properties of that formulation are at odds
with the sets of Section~\ref{sec:method:graph}. Solvers certify a binary as
integral only up to a tolerance $\tau$, so a variable reported as $z_{ij}=1$
may be $1-\tau$ and relax every row of its set by $M\tau$; membership is
meaningful only for sets much thicker than that, whereas transition sets are
deliberately thin. And the continuous relaxation weakens as $M$ grows, so the
search tree grows with it. Tight per-row constants and extended formulations
that model the union by its convex hull both help~\cite{vielma2015mixed}, at
the price of a larger model, but neither removes the outer loop over $\ell$.
The reachability A* has no relaxation and no integrality tolerance, and the
number of grasps is an outcome of the search rather than an input.
